\chapter{Analysis and Theoretical Foundation}
\label{ch:analysis}
\pagestyle{fancy}

\section{General Considerations}

\subsection{Representation and Coordinate Systems} \label{analysis:representation}

Before we deep dive into anything, it is important to lay down some conventions. First of all, following the rules of chess, a chessboard should be positioned in a way that its upper left corner is white. This applies to both players (if we get one right, the other one just stays on the other side). There are different kinds of chess notations to refer to a movement. Usually, they all obey the convention that the rows are numbered, while the columns are letters, and the white player is used as a reference point. In Figure \ref{fig:representation} we can imagine this as the white player being at the ``bottom part'' of the board. However, for us, from a programmer perspective, looking at a chessboard like a matrix, it makes more sense to use numbers in ``both directions''. Since the project has OpenCV at its core, a Mat-like indexation would refer to \t{A8} as \t{00}, \t{E4} as \t{44}, \t{H2} as \t{67} and so on.

\begin{figure}[H]
    \centering
    \includegraphics[width=0.55\textwidth]{figs/representation.png}
    \caption{Representation}
    \label{fig:representation}
\end{figure}

The project uses two cameras. These are mounted on the two sides of the board where the players are not playing, for convenience. If we use the white player as a reference, we can refer to them as left and right cameras. Each camera has its own coordinate system. This is shown in figure \ref{fig:coordinates}. The left camera and its coordinate system are red, while the right camera has the color green associated.

\begin{figure}[H]
    \centering
    \includegraphics[width=\textwidth]{figs/coordinates.png}
    \caption{Coordinates}
    \label{fig:coordinates}
\end{figure}

We need a global coordinate system, and it makes sense to use the one we started with. So the coordinates the cameras see are translated to the white player's coordinate system, which we will refer to as the main coordinate system from now on. If we observe the pattern, we realize that this translation is trivial. If we convert from left to main, $mainRow = leftCol, mainCol = 7 - leftRow$. If we go from right to main, $mainRow = 7 - rightCol, mainCol = rightRow$.





\section{Image Processing Algorithms}
\subsection{Hough} \label{analysis:hough}

The process of line identification from a source image is not at all a trivial problem. The article \i{A survey of Hough Transform} \cite{hough_survey} takes a comprehensive look at this algorithm. It describes it as ``an ingenious method that converts such global curve detection problem into an efficient peak detection problem in parameter space''.

Most of the time, the input should be an edge image. In our case, we apply Canny Edge Detection beforehand. The idea is to define a parameter space, map the image in it, and then take each object pixel, and let them vote about on which lane they think they can lie on. Bins, sometimes called accumulators, count these votes, and the lines with the most votes or above a certain threshold win. This way, any parametric analytic curve could be described- there are applications of Hough on finding circles and ellipses. However, for us, the goal is to get the lines defined by the chessboard, to identify it.

For line parametrization, there were attempts to use different equations to describe a line. The simple slope-intercept form $f(x,y) = y - mx - c = 0$ for example suffered from the limitation that lines perpendicular to the abscissa were impossible to represent. There are workarounds to avoid this, however, I chose to follow the recommendation of \cite{prs_labs} and go with a different form, the one recommended by Duda and Hart: the normal parametrization/ $\rho - \theta$ representation:

\[ \rho = x \cos{\theta} + y \sin{\theta}\]

To illustrate how this representation works, here is Figure \ref{fig:rhotheta1}, showing an example. 

\begin{figure}[H]
    \centering
    \includegraphics[width=0.3\textwidth]{figs/rhotheta1.png}
    \caption{Line in parametric representation}
    \label{fig:rhotheta1}
\end{figure}

The origin is on the top left corner, to follow the convention enforced by openCV. You can see the \t{x} and \t{y} axes, and the red line is the one we want to describe. We take the normal to the red line from the origin: this is the purple line. The length of the purple line is $\rho$. The angle between the \t{x} axis and the purple line is $\theta$. In this particular case, if we take one unit on the axis to be of length 1, we can say that $\rho = 4\sqrt{2}$ and $\theta = \pi/2 = 45^{\circ}$.

Now we just need to define a convention for the ranges we want to check. I will use Figure \ref{fig:rhotheta2} to show some examples, including border cases.

\begin{figure}[H]
    \centering
    \includegraphics[width=0.9\textwidth]{figs/rhotheta2.png}
    \caption{Parameter bounaries}
    \label{fig:rhotheta2}
\end{figure}

The coordinate system is the same as in Figure \ref{fig:rhotheta1}. The dotted box shows the region of interest, more specifically, pixels that are inside the image. Lines as we know go to infinity, so they naturally go out of bounds. On the left picture, I wanted to focus on the limits of $\theta$. We can see that lines \t{L3}, \t{L4} and \t{L5} will be inside the image. While lines \t{L1} and \t{L2} could be described by a normal form, they are out of bounds, hence not of interest. If we think about $\theta_1$ and $\theta_2$, we can observe that they are in the range $(2\pi, 3\pi/4)$, or $(180^{\circ}, 270^{\circ})$. Looking at the right figure, we can focus on $\rho$. Again, some lines are out of the image, and we don't care about them. We cannot filter them based only on $\rho$. But we can see the two extremes. The minimum $\rho$ is achieved by \t{L1}: it has only one pixel inside the image, and it coincides with the origin of the coordinate system, so it is at 0 distance from it. Meanwhile, line \t{L5} is at maximum distance, given that its perpendicular coincides with the diagonal of the image. Overall, we can conclude:

\[\theta \in [0, 2\pi] \cup [3\pi/4, 2\pi)\]
\[\rho \in [0,\sqrt{width * width + height * height}]\]

To enhance the efficiency, some use kernel operations to calculate votes. Others try to use a rotation-based approach, where ``the edge image is rotated about the center of the image plane by a fixed angle step, and in each rotated image, search is carried out in the horizontal and vertical directions to find the end points of horizontal and vertical line segments'' \cite{hough_survey}. Some try to use Neural Networks, to train a model to map an image to parametric lines. And as always, dedicated hardware and parallelization can be of great help.

Overall, the Hough line detection algorithm is somewhat intuitive, robust to discontinuity and missing data on the curve/line, versatile and well-researched. As some downsides, we can mention that the computation time and memory requirements grow exponentially with the number of parameters, and in its simplest form, it suffers from peak spreading (votes in an area for ``the same thing'') and blind voting (every vote has the same weight).

\subsection{Canny Edge Detection}
\subsection{Perspective Transform}

\section{Machine Learning/ Classification Algorithms}
\subsection{K Nearest Neighbors}
\subsection{Support Vector Machine}

Hello
\subsection{Convolutional Neural Network}

\section{Chess Move Validation}

\section{Requirements and Features}
\subsection{Functional Requirements and Features}
\subsection{Nonfunctional Requirements and Features}


